\subsection{自动系数 Dirichlet 多项式：高阶矩、大值与双通道大筛输入}\label{subsec:automatic-dirichlet-polynomial-large-sieve}

本小节把“有限状态生成的系数”作为同一机制核的输出，并把其解析侧输入组织为三条命题链：\textbf{高阶矩节省}（抑制大值频次）、\textbf{临界区间结构性改进}（击穿瓶颈尺度），以及\textbf{双通道大筛输入}（同时压制加性与乘性障碍）。
这些结论的共同假设是：相应的扭曲转移算子在排除有限模障碍后具有统一谱隙；在本稿语义中，这等价于“共振来源仅来自有限可枚举的 cyclotomic 因子族”。

\begin{proposition}[自动系数 Dirichlet 多项式的非渐近矩界：统一 \(2k\) 阶控制]\label{prop:conclusion69-automatic-dirichlet-moment}
设 \((b_n)\) 由某个 finite-state 系统生成，并在区间 \(n\sim N\)（例如 \(N\le n<2N\)）上截断。
令
\[
D(t):=\sum_{n\sim N} b_n\,n^{\mathrm{i}t}.
\]
假设对所有相关扭曲（排除一族显式可枚举的有限模障碍）对应的转移算子族存在统一谱隙。
则对任意固定 \(k\in\NN\) 存在常数 \(\varepsilon_k>0\) 与 \(C_k>0\) 使得对所有 \(T\ge 1\) 有非渐近矩界
\[
\boxed{\;
\int_0^T |D(t)|^{2k}\,dt
\le C_k\,T\,N^{k}\,(\log N)^{C_k}\,N^{-\varepsilon_k}.\;}
\]
该结论提供一条“谱隙 \(\Rightarrow\) 统一高阶矩节省”的可审计输入，可用于将大值频次压缩到均方尺度。
\end{proposition}

\begin{proof}
有限状态生成保证 \(b_n\) 可被编码为有限维扭曲算子族的轨道读出；统一谱隙给出对相关频率扭曲的指数衰减，从而在大筛/均值值分解中产生幂次节省。
将谱隙节省沿 \(2k\) 阶展开（或多线性化）传播即可得到所示非渐近矩界。证毕。
\end{proof}
\leanverified{paper_conclusion69_automatic_dirichlet_moment}

\leanverified{paper\_conclusion70\_critical\_interval\_structure\_improvement}
\begin{theorem}[临界区间的结构性改进定理：自动系数下的严格幂次节省]\label{thm:conclusion70-critical-interval-structure-improvement}
在命题 \ref{prop:conclusion69-automatic-dirichlet-moment} 的设定下，
令
\[
W:=\{t\in[0,T]:\ |D(t)|\ge V\}
\]
为 \(1\)-分离的大值点集。
取临界标度 \(T=N^{6/5}\) 与阈值 \(V\asymp N^{3/4}\)。
则存在 \(\epsilon>0\)（仅依赖谱隙数据与 finite-state 系统）使得
\[
\boxed{\;
|W|\ll N^{12/5-\epsilon}\,V^{-4}.\;}
\]
该不等式在“解析大值瓶颈”的临界项上给出严格幂次改进，并把改进强度参数化为统一谱隙常数与有限模障碍的排除。
\end{theorem}

\begin{proof}
以 \(2k\) 阶矩界为输入，结合 \(1\)-分离与局部化分解，将大值点集的计数约束转写为一族可控的多线性平均；统一谱隙在关键平均上提供额外幂次节省，从而在临界尺度 \(T=N^{6/5}\)、阈值 \(V\asymp N^{3/4}\) 处击穿传统上界。证毕。
\end{proof}

\begin{proposition}[双大筛输入的统一形式：加法角色与乘法角色的同时谱隙化]\label{prop:conclusion74-double-large-sieve-input}
设 \((b_n)\) 由同一 finite-state 系统生成。
假设：
\begin{enumerate}
  \item \textbf{（乘性通道）}对所有非平凡 Dirichlet 角色 \(\chi\ (\bmod q)\)（排除一族有限模障碍）相应扭曲转移算子存在统一谱隙；
  \item \textbf{（加性通道）}对所有加性角色 \(e(an/q)=\exp(2\pi\mathrm{i}an/q)\)（同样排除有限障碍族）相应扭曲转移算子存在统一谱隙。
\end{enumerate}
则对任意 \(Q\le N^\theta\)（固定 \(\theta>0\)），存在 \(\delta_1,\delta_2>0\)（仅依赖谱隙数据）使得对所有 \(q\le Q\) 成立一致节省：
\[
\boxed{\;
\sum_{n\sim N} b_n\chi(n)\ll N^{1-\delta_1},\qquad
\sum_{n\sim N} b_n e(an/q)\ll N^{1-\delta_2}\;}
\]
其中 \((a,q)=1\) 且 \(\chi\) 非平凡。
该结论给出可直接喂入三阶迹/仿射平均体系的“双通道大筛”输入，用于同时压制主弧与乘性障碍。
\end{proposition}
\leanverified{paper\_conclusion74\_double\_large\_sieve\_input}

\begin{proof}
两类角色扭曲都可被编码为有限维扭曲算子族（或等价的有限状态同步编译核）。
统一谱隙给出对所有 \(q\le Q\) 的一致指数衰减；将该衰减经 Perron 分支与谱投影传播到相应的指数和/Dirichlet 和，得到 \(N^{1-\delta}\) 型一致节省。证毕。
\end{proof}

\endinput
